\documentclass[10pt,a4paper]{amsart}
\usepackage[margin=19mm]{geometry}
\usepackage{amsmath,amssymb,mathtools}
\usepackage[scr=rsfs,cal=euler]{mathalfa}
\usepackage{xfrac}
\usepackage[unicode,hidelinks,bookmarks=false]{hyperref}
\newcommand{\ie}{i.e.\ }
\newcommand{\cf}{cf.\ }
\newcommand{\resp}{respectively\ }
\newcommand{\Subsection}{\subsection}
\providecommand{\nobreakdash}{\nobreak\hbox{-}}
\setlength{\emergencystretch}{2em}
\multlinegap0pt
\usepackage{tikz}
\usepackage{amssymb}
\usepackage{enumerate}
\usepackage{float}
\usepackage{xcolor}
\newtheorem{thm}{Theorem}
\newtheorem{lem}{Lemma}
\title{Ratio of the number of $\mathbf{1}$-nearly independent vertex subsets and the Merrifield-Simmons index}
\author{Audace A. V. Dossou-Olory, Eric O. D. Andriantiana, Valisoa R. M. Rakotonarivo, Zekhaya B. Shozi}
\date{Source: arXiv:2604.11813v1 (CC BY 4.0)}
\begin{document}
\maketitle

\noindent\emph{Source and licence.} Ratio of the number of $\mathbf{1}$-nearly independent vertex subsets and the Merrifield-Simmons index. Version arXiv:2604.11813v1 (CC BY 4.0), distributed by arXiv under the Creative Commons Attribution 4.0 International licence, http://creativecommons.org/licenses/by/4.0/, as recorded in the arXiv licence metadata for this exact version and shown on the abstract page https://arxiv.org/abs/2604.11813v1. Full licence notice: \texttt{licenses/arxiv-2604.11813-CC-BY-4.0.txt}.\par\smallskip

\noindent\textit{Editorial scope.} Counted unit: the article's thm MaxForest (source0.tex lines 463--469). The article's complete original proof of it is delivered with it (source0.tex lines 522--564, one \texttt{proof} environment of 43 lines, which the article places away from the statement). No mathematics is written, repaired or re-derived: the statement and the proof are the article's own lines, in the article's own notation, and no retained line is substituted or re-flowed.  Retained uncounted as the source-local context the proof needs: lines 474--479; lines 497--502.  Nothing was omitted from any retained span, so the package carries no omission record.  No retained line contains a citation, so the wrapper carries no bibliography and no bibliography entry.  No retained line contains a figure environment, an image-inclusion command or drawing code, so no image is delivered and the package declares no asset dependency.  Editorial formatting only, in addition: an AMS \texttt{amsart} preamble that restates the article's own declarations and macros used by the retained lines, namely the environments \texttt{thm}, \texttt{lem} on separate counters, together with the standard packages those lines need.  The retained environments print in this document as Theorem 1, Lemma 1 in order of appearance, whereas the article prints the same items with the numbering of its own full document.  The complete native source of the article as distributed in the arXiv e-print for this exact version is bundled unchanged as \texttt{sources/198362/source0.tex}.
	\begin{lem}\label{Lem:1}
		If $T$ is a tree of order $n\geq 2$  and $v$ a leaf of $T$, then it holds that
		\begin{equation*}
			\frac{\sigma_0\left(T-N_T[v]\right)}{\sigma_0\left(T-v\right)}\leq 1-\frac{1}{2^{n-2}+1}\,.
		\end{equation*}
	\end{lem}

	\begin{lem}\label{Lem:2}
		If $T$ is a tree and $v$ a leaf of $T$, then it holds that
		\begin{equation*}
			Q(T)\leq \frac{\displaystyle\frac{\sigma_0\left(T-v\right)}{\sigma_0\left(T-N_T[v]\right)}Q(T-v)+1+Q\left(T-N_T[v]\right)}{1+\displaystyle\frac{\sigma_0\left(T-v\right)}{\sigma_0\left(T-N_T[v]\right)}}\,.
		\end{equation*}
	\end{lem}

	\begin{thm}\label{MaxForest}
		Let $F$ be a forest of order $n\geq 1$. We have
        \begin{equation*}
			Q(F)\leq \frac{1}{3}\left(n-1\right)\,.
		\end{equation*}
The equality is reached for $F\in \{P_1,P_2\}$.
	\end{thm}

	\begin{proof} [Proof of Theorem~\ref{MaxForest}]
		We proceed by induction on the order of the forest $F$. Let $n=|F|:=|V(F)|$. If $n=1$, then $F=P_1$, $Q(F)=Q(P_1)=0$, and $0= \frac{1}{3}(1-1)$. If $n=2$, then $F=P_2$ or $F=P_1\cup P_1$. We have
			$Q(P_2)=\frac{1}{3}$, while $\frac{1}{3}\left(2-1\right)=\frac{1}{3}$. Also
			$Q(P_1\cup P_1)=0<\frac{1}{3}\left(2-1\right)$.
            
			Assume that the theorem holds for all forests with order less than or equal to $n-1$, for some $n\geq 3$. Let $F$ be a forest of order $n$. If $F$ is not connected , then $F=\cup_{i=1}^mT_i$, for some $m\geq 2$. In this case,
			\begin{equation*}
			Q(F)=\sum_{i=1}^{m}Q(T_i)\leq \sum_{i=1}^{m}\frac{1}{3}\left(|T_i|-1\right)
					=\frac{1}{3}|T|-\frac{1}{3}m < \frac{1}{3}|T|-\frac{1}{3}\,.
			\end{equation*}
			So, assume $F$ is a tree $T$. 
    %\medskip
	%We continue the proof of the theorem.
    Set $a:=\frac{1}{3}$. By Lemma~\ref{Lem:2} together with the induction hypothesis, we get
		\begin{equation*}
	\begin{aligned}
			Q(T)&\leq \frac{\displaystyle\frac{\sigma_0\left(T-v\right)}{\sigma_0\left(T-N_T[v]\right)}Q(T-v)+1+Q\left(T-N_T[v]\right)}{1+\displaystyle\frac{\sigma_0\left(T-v\right)}{\sigma_0\left(T-N_T[v]\right)}}\\
			&\leq \frac{\displaystyle\frac{\sigma_0\left(T-v\right)}{\sigma_0\left(T-N_T[v]\right)}(a(n-2))+1+a(n-3)}{1+\displaystyle\frac{\sigma_0\left(T-v\right)}{\sigma_0\left(T-N_T[v]\right)}}\,.
	\end{aligned}
	\end{equation*}
    
	Using Lemma~\ref{Lem:1}, that is $$\frac{\sigma_0\left(T-v\right)}{\sigma_0\left(T-N_T[v]\right)}\geq \frac{1}{\displaystyle 1-\frac{1}{2^{n-2}+1}}=\frac{2^{n-2}+1}{2^{n-2}}\,,$$ 
    together with the fact that the function $x \mapsto (Ax+B)/(1+x)$ is decreasing  if $A< B$, we obtain
	\begin{equation*}
		Q(T)\leq \frac{\displaystyle\frac{2^{n-2}+1}{2^{n-2}}\left(a(n-2)\right)+ 1+a(n-3)}{1+\displaystyle \frac{2^{n-2}+1}{2^{n-2}}}\,.
	\end{equation*}
	Set $c_n:=(2^{n-2}+1) / 2^{n-2}$. We have
	\begin{equation*}
		Q(T)\leq \frac{c_n\left(an-a-a\right)+an-2a-a+1}{1+c_n}.
	\end{equation*}
  We finish the proof by showing that  
$$c_n\left(an-a-a\right)+an-2a-a+1\leq \left(1+c_n\right)\left(an-a\right)\,.$$
This simplifies to 
%This is equivalent to showing that
	\begin{equation*}
		\begin{aligned}
			%& c_n(-a)+an-2a-a+1\leq an-a\,,\\
			& ac_n+2a-1\geq 0\,.
		\end{aligned}
	\end{equation*}
	Since $c_n\geq 1$ for all positive integers $n$, it suffices to show that $a(1)+2a-1 \geq 0$, that is $a\geq \frac{1}{3}$,
	which is indeed true.
	\end{proof}

\end{document}
